%!TEX root = ../nauchka.tex

\section{Вертикальное масштабирование}

%%%%%%%%%%%%%%%%%%%%%%%%%%%%%%%%%%%%%%%%%%%%%%%%%%%%%%%%%%%%%%%%%%%%%%%%%%%%%%%%
\subsection{Оптимизации на уровне одного GPU}

\begin{frame}
  \frametitle{Вертикальное масштабирование: проблема памяти}
  \begin{itemize}
    \item Цель: уменьшить потребление памяти на \textbf{одном GPU}, чтобы поместить бóльшую модель или батч.
    \item Основные потребители памяти при обучении:
      \begin{itemize}
        \item Веса модели: $M_{\text{weights}}$
        \item Градиенты: $M_{\text{grad}}$ (столько же, сколько веса)
        \item Состояния оптимизатора: $M_{\text{opt}}$ (обычно $2 \times M_{\text{weights}}$ для Adam)
        \item Промежуточные активации: $M_{\text{act}}$
          \[
            M_{\text{act}} \propto B \cdot L \cdot S \cdot H 
          \]
          ($B$ - батч, $L$ - число слоев, $S$ - длина последовательности, $H$ - скрытый размер)
          \textbf{Часто являются основным узким местом.}
      \end{itemize}
    \item Методы вертикального масштабирования направлены на сокращение одного или нескольких из этих компонентов.
  \end{itemize}
\end{frame}

%%%%%%%%%%%%%%%%%%%%%%%%%%%%%%%%%%%%%%%%%%%%%%%%%%%%%%%%%%%%%%%%%%%%%%%%%%%%%%%%
\begin{frame}
  \frametitle{Activation Checkpointing (Gradient Checkpointing)}
  \begin{itemize}
    \item Идея: не хранить все активации для обратного прохода.
    \item Вместо этого, некоторые активации \textbf{перевычисляются} во время обратного прохода.
    \item Преимущества:
      \begin{itemize}
        \item Значительно сокращает $M_{\text{act}}$.
        \item Позволяет использовать бóльшие батчи или более длинные последовательности.
      \end{itemize}
    \item Недостатки:
      \begin{itemize}
        \item Увеличивает время обучения из-за дополнительных прямых проходов (обычно ~30\% для трансформеров).
      \end{itemize}
    \item Пример: при чекпоинтинге каждого слоя трансформера, потребление памяти под активации снижается с $O(L)$ до $O(\sqrt{L})$ или даже $O(1)$ в зависимости от стратегии.
  \end{itemize}
  \vspace{0.5cm}
  \begin{flushright}
    \tiny Источник: \url{https://arxiv.org/abs/1911.13214}
  \end{flushright}
\end{frame}

%%%%%%%%%%%%%%%%%%%%%%%%%%%%%%%%%%%%%%%%%%%%%%%%%%%%%%%%%%%%%%%%%%%%%%%%%%%%%%%%
\begin{frame}
  \frametitle{Fusing Operations (Kernel Fusion)}
  \begin{itemize}
    \item Идея: объединить несколько последовательных операций GPU (ядер) в одно «супер-ядро».
    \item Примеры:
      \begin{itemize}
        \item Bias + ReLU $\rightarrow$ FusedBiasReLU
        \item Несколько поэлементных операций.
        \item Слой LayerNorm, Attention-слои могут быть частично слиты.
      \end{itemize}
    \item Преимущества:
      \begin{itemize}
        \item \textbf{Уменьшает нагрузку на память (memory bandwidth):} промежуточные результаты остаются в кэше/регистрах, а не записываются/читаются из глобальной памяти GPU.
        \item \textbf{Снижает накладные расходы на запуск ядер.}
        \item Ускоряет выполнение.
      \end{itemize}
    \item Может незначительно уменьшить пиковое потребление памяти за счет отсутствия промежуточных тензоров в глобальной памяти.
  \end{itemize}
  \vspace{0.5cm}
  \begin{flushright}
    \tiny Источник: \url{https://hazyresearch.stanford.edu/blog/2024-05-12-tk}
  \end{flushright}
\end{frame}

%%%%%%%%%%%%%%%%%%%%%%%%%%%%%%%%%%%%%%%%%%%%%%%%%%%%%%%%%%%%%%%%%%%%%%%%%%%%%%%%
\begin{frame}
  \frametitle{FlashAttention}
  \begin{itemize}
    \item \textbf{I/O-aware exact attention algorithm}, не аппроксимация.
    \item Проблема стандартной реализации attention:
      \begin{itemize}
        \item Требует материализации матрицы внимания $P = \text{softmax}(\frac{QK^T}{\sqrt{d_k}})$ размером $S \times S$.
        \item Потребление памяти для $P$: $O(S^2)$.
        \item Множественные чтения/записи $Q, K, V$ из медленной HBM.
      \end{itemize}
    \item FlashAttention использует \textbf{tiling (разбиение на блоки)} и \textbf{recomputation (перевычисление)}:
      \begin{itemize}
        \item Загружает блоки $Q, K, V$ в быструю SRAM.
        \item Вычисляет блоки матрицы внимания и агрегирует результат на лету, \textbf{не материализуя полную матрицу $P$ в HBM}.
      \end{itemize}
    \item Преимущества:
      \begin{itemize}
        \item Снижает потребление памяти для attention с $O(S^2)$ до $O(S)$ (для промежуточных вычислений, не для хранения активаций $Q,K,V$ для backward pass).
        \item Значительно ускоряет вычисление attention (до ~2-4x).
      \end{itemize}
  \end{itemize}
  \vspace{0.5cm}
  \begin{flushright}
    \tiny Источник: \url{https://arxiv.org/abs/2205.14135}
  \end{flushright}
\end{frame}

% Можно добавить слайд с картинкой для FlashAttention, если найдете подходящую
% \begin{frame}{FlashAttention: Иллюстрация}
%   \begin{figure}
%     \centering
%     \includegraphics[width=0.7\textwidth]{flash_attention_concept.png} % Пример имени файла
%     \caption{Концепция FlashAttention (Tiling)}
%     \label{fig:flash_attention}
%   \end{figure}
%   \vspace{0.5cm}
%   \begin{flushright}
%     \tiny Источник: \url{https://arxiv.org/abs/2205.14135}
%   \end{flushright}
% \end{frame}

%%%%%%%%%%%%%%%%%%%%%%%%%%%%%%%%%%%%%%%%%%%%%%%%%%%%%%%%%%%%%%%%%%%%%%%%%%%%%%%%
\begin{frame}
  \frametitle{Mixed Precision Training}
  \begin{itemize}
    \item Идея: использовать форматы с пониженной точностью (FP16, BF16) для большинства вычислений и хранения, сохраняя FP32 для критических частей.
    \item Компоненты:
      \begin{itemize}
        \item \textbf{FP16/BF16}: веса, активации, градиенты.
        \item \textbf{FP32}: мастер-копия весов (для аккумулирования градиентов и обновлений), вычисление лосса (иногда).
      \end{itemize}
    \item \textbf{Loss Scaling} (особенно для FP16):
      \begin{itemize}
        \item Градиенты в FP16 могут быть слишком малы (underflow).
        \item Лосс умножается на скаляр (scaling factor) перед backward pass.
        \item Градиенты делятся на тот же скаляр перед обновлением весов.
      \end{itemize}
    \item Преимущества:
      \begin{itemize}
        \item В ~2 раза уменьшает $M_{\text{weights}}$, $M_{\text{grad}}$, $M_{\text{act}}$.
        \item Ускоряет вычисления на GPU с Tensor Cores.
      \end{itemize}
  \end{itemize}
  \vspace{0.5cm}
  \begin{flushright}
    \tiny Источник: \url{https://developer.nvidia.com/blog/video-mixed-precision-techniques-tensor-cores-deep-learning/}
  \end{flushright}
\end{frame}

%%%%%%%%%%%%%%%%%%%%%%%%%%%%%%%%%%%%%%%%%%%%%%%%%%%%%%%%%%%%%%%%%%%%%%%%%%%%%%%%
\begin{frame}
  \frametitle{BF16 и FP8 форматы}
  \begin{itemize}
    \item \textbf{BFloat16 (BF16)}:
      \begin{itemize}
        \item Формат: 1 знаковый бит, 8 бит экспоненты, 7 бит мантиссы.
        \item Такой же диапазон экспоненты, как у FP32 $\implies$ более устойчив к переполнению/исчезновению по сравнению с FP16.
        \item Меньшая точность мантиссы, чем у FP16.
        \item Часто не требует loss scaling. Широко используется для обучения LLM.
      \end{itemize}
    \item \textbf{FP8 (8-bit Floating Point)}:
      \begin{itemize}
        \item Несколько вариантов: E4M3 (4 бит экспоненты, 3 мантиссы) и E5M2 (5 бит экспоненты, 2 мантиссы).
        \item E4M3: бóльшая точность, меньший диапазон.
        \item E5M2: меньшая точность, бóльший диапазон.
        \item Требует специальных техник для поддержания качества модели (например, динамическое масштабирование, гибридное использование с BF16).
        \item Обеспечивает еще большее сокращение памяти (~4x по сравнению с FP32) и ускорение.
        \item Поддерживается на новых GPU (Hopper H100+).
      \end{itemize}
  \end{itemize}
\end{frame}
